\documentclass[12pt,a4paper]{article}
\usepackage{graphicx}
\usepackage{amsmath}
\usepackage{xcolor}
\usepackage{geometry}
\usepackage{fancyhdr}
\usepackage{tikz}
\usepackage{enumitem}

\geometry{a4paper, margin=1in}

% Set up the header and footer
\pagestyle{fancy}
\fancyhf{}
\fancyhead[L]{\leftmark}
\fancyhead[R]{\thepage}

\begin{document}
	
	% Title Page
	\begin{titlepage}
		\centering
		\vspace*{2cm}
		
		% Decorative Line
		\noindent\rule{\textwidth}{1pt}\\[0.5cm]
		
		% Project Title
		{\Huge\textbf{\color{blue}LAB NOTEBOOK USING LATEX}}\\
		\vspace{0.5cm}
		\noindent\rule{\textwidth}{1pt}\\[2cm]
		
		% Group Details
		\Large\textbf{Submitted by:}\\
		\vspace{0.5cm}
		
		\textbf{\color{red}Group Leader:}\\
		\vspace{0.3cm}
		\Large MD Erfanuddin\\
		\large Roll Number: 30001223001\\
		Department: BCA\\
		
		\vspace{1.5cm}
		
		\textbf{\color{green}Group Members:}\\
		\vspace{0.5cm}
		
		\large
		\begin{tabbing}
			\hspace{4cm} \= \hspace{8cm} \= \kill
			\textbf{Name} \> \textbf{Roll Number} \> \textbf{Department} \\
			\textcolor{purple}{Sanchita Laha} \> 30084323022 \> BSc in Data Science \\
			\textcolor{purple}{Arnov Pal} \> 30001223067 \> BCA \\
			\textcolor{purple}{Abir Dey} \> 30001223010 \> BCA \\
			\textcolor{purple}{Ibnoon Qais} \> 30059223003 \> BSc in Forensic Sciences \\
		\end{tabbing}
		
		\vfill
		
		% Department and University Information
		\Large\textbf{\color{orange}Git Assignment}\\
		\Large\textbf{\color{blue}Maulana Abul Kalam Azad University of Technology}\\
		\vspace{0.5cm}
		\Large\textbf{Date: \today}
		
		% Decorative Line at the Bottom
		\vspace{2cm}
		\noindent\rule{\textwidth}{1pt}
		
	\end{titlepage}
	
	% Index Page
	\newpage
	\tableofcontents
	
	% Md Erfanuddin's Page
	\newpage
	\section{C Programming Calculator /by : Md.Erfanuddin}
	\subsection{Introduction}
	This document outlines how to set up a local repository, create a C calculator program, commit the code, and push it to GitHub.
	\subsection{Steps}
	\begin{enumerate}
		\item Create a local repository using git init.
		\item Build the C program (Calculator).
		\item Commit the code with git commit.
		\item Create a public repository on GitHub.
		\item Push the local repository to GitHub.
	\end{enumerate}
	\subsection{Code Snippet}
	\begin{verbatim}
		#include <stdio.h>
		
		int main() {
			char operator;
			double num1, num2;
			
			printf("Enter an operator (+, -, *, /): ");
			scanf("%c", &operator);
			printf("Enter two operands: ");
			scanf("%lf %lf", &num1, &num2);
			
			switch (operator) {
				case '+': printf("%.2lf + %.2lf = %.2lf\n", num1, num2, num1 + num2); break;
				case '-': printf("%.2lf - %.2lf = %.2lf\n", num1, num2, num1 - num2); break;
				case '*': printf("%.2lf * %.2lf = %.2lf\n", num1, num2, num1 * num2); break;
				case '/': 
				if (num2 != 0.0) printf("%.2lf / %.2lf = %.2lf\n", num1, num2, num1 / num2);
				else printf("Division by zero error.\n");
				break;
				default: printf("Invalid operator.\n"); break;
			}
			return 0;
		}
		
	\end{verbatim}
% Abir Dey's Page
	\newpage
	\section{Enumerate ABC Format and Roman Number / by: Abir Dey}
	\subsection{Introduction}
	This document demonstrates how to create enumerated lists in LaTeX using ABC format and Roman numerals, providing flexibility in list formatting.
	
	\subsection{Creating an ABC Format Enumerated List}
	To create an enumerated list in ABC format (a, b, c), use the enumerate environment with label=\alph*.:
	
	\begin{verbatim}
		\begin{enumerate}[label=\alph*.]
			\item First item
			\item Second item
			\item Third item
		\end{enumerate}
	\end{verbatim}
	
	Result:
	
	\begin{enumerate}[label=\alph*.]
		\item First item
		\item Second item
		\item Third item
	\end{enumerate}
	
	\subsection{Creating a Roman Numerals Enumerated List}
	For a list with Roman numerals (i, ii, iii), use label=\roman*.:
	
	\begin{verbatim}
		\begin{enumerate}[label=\roman*.]
			\item First item
			\item Second item
			\item Third item
		\end{enumerate}
	\end{verbatim}
	
	Result:
	
	\begin{enumerate}[label=\roman*.]
		\item First item
		\item Second item
		\item Third item
	\end{enumerate}
 % Arnov Pal's Page
	\newpage
	\section{Other Styles Such as Square and Circle /by: Arnov pal}
	\subsection{Introduction}
	This document demonstrates the creation of simple geometric shapes, such as squares and circles, using the TikZ package in LaTeX. Such shapes are often used in diagrams, flowcharts, and illustrations.
	
	\subsection{Examples}
	\subsubsection{Square}
	\begin{center}
		\begin{tikzpicture}
			\draw[fill=blue!30, thick] (0,0) rectangle (3,3);
		\end{tikzpicture}
	\end{center}
	
	\subsubsection{Circle}
	\begin{center}
		\begin{tikzpicture}
			\draw[fill=red!30, thick] (0,0) circle (1.5cm);
		\end{tikzpicture}
	\end{center}
	
	\subsection{Conclusion}
	The ability to create and customize shapes like squares and circles in LaTeX using TikZ provides great flexibility for creating professional documents and presentations.
 % Sanchita Laha's Page
	\newpage
	\section{Create 5 Subsections/ by : Sanchita Laha}
	\subsection{Introduction}
	
	\subsubsection{Subsection 1: Overview}
	In this subsection, provide an overview of the topic you're discussing. This is the first of the five subsections.
	
	\subsubsection{Subsection 2: Background Information}
	Here, you can include any necessary background information that the reader needs to understand the topic.
	
	\subsubsection{Subsection 3: Methodology}
	Describe the methods or procedures used in your work or research. This is where you detail how you approached the problem.
	
	\subsubsection{Subsection 4: Analysis}
	Present the analysis of the data or information you've gathered. Discuss any patterns, trends, or insights you discovered.
	
	\subsubsection{Subsection 5: Conclusion}
	Summarize your findings and discuss the implications of your work. Mention any future directions or next steps.
    \end{document}
